\BourbakiExercisesSection

\begin{enumerate}
\def\labelenumi{\arabic{enumi}.}
\tightlist
\item
  \begin{enumerate}
  \def\labelenumii{(\alph{enumii})}
  \tightlist
  \item
    设 E 是一个右 A-模。在加法群 \(E^{r}\) （ \(r \geqslant 1\) ）上定义一个外合成法则，以 \(\mathbf{M}_{r}(\mathbf{A})\) 为算子集：对每个元素 \(x = (x_{i})_{1 \leqslant i \leqslant r}\) （属于 \(E^{r}\) ）和每个方阵 \(P = (\alpha_{ij}) \in \mathbf{M}_{r}(\mathbf{A})\) ，用 x.P 记元素 \(y = (y_{i})\) （属于 \(E^{r}\) ），使得
  \end{enumerate}
\end{enumerate}

\[
y _ {i} = \sum_ {j = 1} ^ {r} x _ {j} \alpha_ {j i} \quad (1 \leqslant i \leqslant r).
\]

该外合成法则与 \(E^{r}\) 上的加法法则一起定义一个右 \(\mathbf{M}_{r}(\mathbf{A})\) -模结构（在 \(E^{r}\) 上）；把算子环限制为标量矩阵 \(I.\alpha\) 所成的环，就重新得到 \(E^{r}\) 上的积 A-模结构。证明：为使 A-模 E 有一个 r 个元素的生成系，必要且充分条件是 \(\mathbf{M}_{r}(\mathbf{A})\) -模 \(E^{r}\) 是单生成的。

\begin{enumerate}
\def\labelenumi{(\alph{enumi})}
\setcounter{enumi}{1}
\tightlist
\item
  设 \((x_{i})_{1\leqslant i\leqslant n}\) 是 E 的一个生成系，而 \((y_{j})_{1\leqslant j\leqslant m}\) 是 E 的一个元素族（相应地，E 的一个生成系）。设 \(z, z', z''\) 是 \(\mathbf{E}^{m + n}\) 的三个元素，使得 \(z_{i} = x_{i}\) 对 \(1\leqslant i\leqslant n\) , \(z_{n + j} = 0\) 对 \(1\leqslant j\leqslant m\) , \(z_{i}' = x_{i}\) 对
\end{enumerate}

\(1 \leqslant i \leqslant n, z_{n+j}^{\prime} = y_{j}\) 对 \(1 \leqslant j \leqslant n, z_{i}^{\prime \prime} = 0\) 对 \(1 \leqslant i \leqslant n, z_{n+j}^{\prime \prime} = y_{j}\) 对 \(1 \leqslant j \leqslant n\) 。证明存在两个可逆矩阵 \(P, Q\) （属于 \(\mathbf{M}_{n+m}(\mathbf{A})\) ），使得 \(z' = z.P\) （相应地， \(z'' = z.Q\) ）。

\begin{enumerate}
\def\labelenumi{(\alph{enumi})}
\setcounter{enumi}{2}
\tightlist
\item
  若 A 交换，则 \(u_{P}: x \mapsto x \cdot P\) 是 A-模 \(E^{r}\) 的一个自同态，且映射 \(P \mapsto u_{P}\) 是环 \(\mathbf{M}_{r}(\mathbf{A})\) 到环 \(\operatorname{End}_{\mathbf{A}}(\mathbf{E}^{r})\) 的一个同态。若 E 是忠实 A-模，则该同态是单射。
\end{enumerate}

\begin{enumerate}
\def\labelenumi{\arabic{enumi}.}
\setcounter{enumi}{1}
\tightlist
\item
  \begin{enumerate}
  \def\labelenumii{(\alph{enumii})}
  \tightlist
  \item
    设 \(X\) 是环 \(\mathbf{A}\) 上的一个方阵，它可以写成一个上三角矩阵 \((X_{ij})\) （由矩阵组成 \((1 \leqslant i \leqslant p, 1 \leqslant j \leqslant p)\) ）。证明：若每个方阵 \(X_{ii}\) \((1 \leqslant i \leqslant p)\) 都可逆，则 \(X\) 也可逆，且 \(X^{-1}\) 可写成一个上三角矩阵 \((Y_{ij})\) \((1 \leqslant i \leqslant p, 1 \leqslant j \leqslant p)\) ，对应于指标集的同一分划。当 \(\mathbf{A}\) 是域时，证明这一使 \(X\) 可逆的充分条件也是必要的。
  \end{enumerate}
\end{enumerate}

\begin{enumerate}
\def\labelenumi{(\alph{enumi})}
\setcounter{enumi}{1}
\tightlist
\item
  给出一个环 A 和一个矩阵 \(\begin{pmatrix} a & 1 \\ 0 & b \end{pmatrix}\) （其中 \(a \in \mathbf{A}, b \in \mathbf{A}\) ），该矩阵可逆，但 \(a\) 或 \(b\) 在 A 中均不可逆，且其逆矩阵 \(\begin{pmatrix} a' & b' \\ c' & d' \end{pmatrix}\) 满足 \(b' \neq 0\) 和 \(c' \neq 0\) （取 A 为无限维向量空间的自同态环）。
\end{enumerate}

\begin{enumerate}
\def\labelenumi{\arabic{enumi}.}
\setcounter{enumi}{2}
\tightlist
\item
  设 A 为商环 Z/30Z。证明在矩阵
\end{enumerate}

\[
\left( \begin{array}{c c c} 1 & 1 & - 1 \\ 0 & 2 & 3 \end{array} \right)
\]

上，两行线性无关，但任意两列线性相关。

\begin{enumerate}
\def\labelenumi{\arabic{enumi}.}
\setcounter{enumi}{3}
\tightlist
\item
  设 A 为一个环，C 为其中心，B 为矩阵环 \(\mathbf{M}_n(\mathbf{A}), \Delta\) 为由元素 \(\alpha\beta - \beta\alpha\) （对 \(\alpha \in \mathbf{A}, \beta \in \mathbf{A}\) ）生成的 A 的加法子群，D 为由矩阵 \(XY - YX\) （对 \(X \in \mathbf{B}, Y \in \mathbf{B}\) ）生成的 B 的加法子群； \(\Delta\) 和 D 均为 C-模。对每个矩阵 \(X = (\xi_{ij}) \in \mathbf{B}\) ，设 \(\theta(X)\) 为元素 \(\sum_{i=1}^{n} \bar{\xi}_{ij}\) （属于 A/ \(\Delta\) ），其中对所有 \(\alpha \in \mathbf{A}\) , \(\bar{\alpha}\) 表示 \(\alpha\) 模 \(\Delta\) 的类。证明 \(\theta\) 是 B 到 A/ \(\Delta\) 上的一个满同态，其核等于 D，因而 B/D 作为 C-模同构于 A/ \(\Delta\) 。（注意 D 包含矩阵 \(\alpha E_{ij}\) 以及矩阵 \(\alpha(E_{ii} - E_{jj})\) ，对 \(\alpha \in \mathbf{A}\) 和 \(i \neq j\) 。）
\end{enumerate}

¶ 5. 设 A 是一个环，L、M、N 为任意三个指标集， \(U = (a_{\lambda \mu})\) 是 \(\mathbf{A}^{\mathbf{L} \times \mathbf{M}}\) 中的一个矩阵， \(V = (b_{\mu \nu})\) 是 \(\mathbf{A}^{\mathbf{M} \times \mathbf{N}}\) 中的一个矩阵；若对每个有序对 \((\lambda, \nu) \in \mathbf{L} \times \mathbf{N}\) ，族 \((a_{\lambda \mu} b_{\mu \nu})\) （ \(\mu \in \mathbf{M}\) ）有有限支撑，则元素 \(c_{\lambda \nu} = \sum_{\mu} a_{\lambda \mu} b_{\mu \nu}\) 有定义，这时也称矩阵 \((c_{\lambda \nu})\) 是乘积 \(UV\) —— \(U\) 乘以 \(V\) 的乘积。当乘积 \(UV'\) 和 \(UV''\) 都有定义时， \(U(V' + V'')\)

和 \(UV' + UV'' = U(V' + V'')\) ；逆命题是否成立？给出三个无限矩阵 U、V、W 的例子，使得乘积 UV、VW、 \(U(VW)\) 和 \((UV)W\) 都有定义，但 \(U(VW) \neq (UV)W\) （取 U 为所有元素都等于 1 的单行矩阵，W 为其转置，V 为所有元素都是 0、1 或 -1 且每行每列只有有限个元素 \(\neq 0\) 的矩阵。用归纳法确定 V 的元素，使得 UV = 0 而 VW 只有一个元素 \(\neq 0\) 。）

\begin{enumerate}
\def\labelenumi{\arabic{enumi}.}
\setcounter{enumi}{5}
\item
  设 X 是域 K 上一个 m 行 n 列的矩阵；证明秩 \(\operatorname{rg}(X)\) 等于 X 的行数与列数相同的子矩阵的秩中的最大者。（设 \(\operatorname{rg}(X) = r\) ；若 \(a_{1}, \ldots, a_{r}\) 是 X 的 r 个列，在 \(K_{d}^{m}\) 中组成一个自由系，且用这 r 个向量和典范基的 m - r 个向量构成 \(K_{d}^{m}\) 的一组基，证明 \(a_{1}, \ldots, a_{r}\) 在典范基的其余 r 个向量上的分量构成一个秩为 r 的矩阵。）
\item
  设 X 是域 K 上的一个 m 行 n 列矩阵；若 r 是 X 的秩，证明通过删去 X 的 n - s 列得到的 m 行 s 列子矩阵的秩 \(\geqslant r + s - n\) 。
\item
  设 \(X = (\alpha_{ij})\) 是一个 \(m\) 行 \(n\) 列的矩阵（在域 \(K\) 上）。为使 \(X\) 的秩为 1，必要且充分条件是：在 \(K\) 中存在一个族 \((\lambda_i)_{1 \leqslant i \leqslant m}\) （由 \(m\) 个不全为零的元素组成）和一个族 \((\mu_j)_{1 \leqslant j \leqslant n}\) （由 \(n\) 个不全为零的元素组成），使得对每对有序指标有 \(\alpha_{ij} = \lambda_i \mu_j\) 。
\item
  设 X、Y 是域 K 上 m 行 n 列的两个矩阵；若存在两个 m 阶方阵 \(P, P_{1}\) 和两个 n 阶方阵 \(Q, Q_{1}\) ，使得 Y = PXQ 且 \(X = P_{1}XQ_{1}\) ，证明 X 与 Y 等价。
\item
  设 \(X, X', Y, Y'\) 是四个 \(n\) 阶方阵（在环 \(A\) 上），使得 \(X\) 可逆。为使存在两个可逆方阵 \(P, Q\) （ \(n\) 阶）满足 \(X' = PXQ\) 和 \(Y' = PYQ\) ，必要且充分条件是 \(X'\) 可逆且矩阵 \(YX^{-1}\) 与 \(Y'X'^{-1}\) 相似。
\item
  设 E 是域 K 上维数 \(\geqslant1\) 的一个右向量空间，H 是 E 的一个超平面。E 的每个使 H 的每个元素都不变的自同态 u 在过渡到商时给出 1 维商空间 E/H 的一个自同态，因而它具有形式 \(\dot{x}\mapsto\dot{x}\mu(\dot{x})\) ，其中 \(\mu(\dot{x})\in\mathrm{K}\) 满足 \(\mu(\dot{x}\lambda)=\lambda^{-1}\mu(\dot{x})\lambda\) （对 \(\lambda\in K^{*}\) ）。若 E 的一个使 H 的元素都不变的自同构所对应的 E/H 的自同构是恒等映射，则称它为超平面 H 的平延；否则称为超平面 H 的伸缩。当 u 是一个伸缩时，元素 \(\mu(\dot{x})\) （对 \(\dot{x}\in E/H\) ）所成的集——它是乘法群 \(K^{*}\) 中共轭元素的一个类——称为伸缩 u 的类。
\end{enumerate}

\begin{enumerate}
\def\labelenumi{(\alph{enumi})}
\item
  证明对每个伸缩，存在一条且仅一条 H 的补线，它在该伸缩下不变。
\item
  设 \(\phi\) 是 \(\mathbf{E}\) 上满足 \(\mathbf{H} = \overline{\phi}(0)\) 的一个线性形式；证明对每个平延 \(u\) （超平面 \(\mathbf{H}\) 的平延），存在唯一的向量 \(a \in \mathbf{H}\) ，使得 \(u(x) = x + a\phi(x)\) （对所有 \(x \in \mathbf{E}\) ）。设 \(\Gamma(\mathbf{E}, \mathbf{H})\) 是 \(\mathbf{GL}(\mathbf{E})\) 的由使 \(\mathbf{H}\) 的每个元素不变的自同构组成的子群；证明超平面 \(\mathbf{H}\) 的平延组成一个正规交换子群 \(\Theta(\mathbf{E}, \mathbf{H})\) （ \(\Gamma(\mathbf{E}, \mathbf{H})\) 的子群），它同构于加法群 \(\mathbf{H}\) ；商群 \(\Gamma(\mathbf{E}, \mathbf{H}) / \Theta(\mathbf{E}, \mathbf{H})\) 同构于乘法群 \(\mathbf{K}^*\) 。为使 \(\Gamma(\mathbf{E}, \mathbf{H})\) 交换，必要且充分条件是下列两个条件之一成立：（ \(\alpha\) ) \(\mathbf{K}\) 是含两个元素的域；（ \(\beta\) ) \(\mathbf{K}\) 交换且 \(\dim \mathbf{E} = 1\) 。
\item
  设 E 是有限维的；证明对每个平延 u，存在 E 的一组基，使得 u 在这组基下的矩阵的所有对角元素都等于 1，且至多还有一个别的元素 \(\neq0\) 。
\item
  证明群 \(\Theta(E, H)\) 在群 \(\mathbf{GL}(E)\) 中的中心化子（I，§ 5，第 3 号）是复合 \(Z(E)\,\Theta(E, H) = \Theta(E, H)Z(E)\) —— \(\Theta(E, H)\) 与中心 \(Z(E)\) （ \(\mathbf{GL}(E)\) 的中心）的复合（参见 § 1，习题 26）。属于该中心化子且使至少一个元素 \(\neq0\) （属于 \(E\) ）不变的自同构只有群 \(\Theta(E, H)\) 的平延。若 \(K\) 至少有 3 个元素，则 \(\Gamma(E, H)\) 在 \(\mathbf{GL}(E)\) 中的中心化子等于 \(Z(E)\) 。
\item
  证明 \(\Theta(\mathbf{E},\mathbf{H})\) 和 \(\Gamma(\mathbf{E},\mathbf{H})\) 在 \(\mathbf{GL}(\mathbf{E})\) 中的正规化子（I，§ 5，第 3 号）两者都等于由使 H 不变的自同构组成的子群。
\end{enumerate}

¶ 12. 设 E 是域 K 上维数 ≥1 的一个右向量空间。用 F(E) 记 GL(E) 的由使 1\_E - u 的核具有有限余维数的自同构 u 组成的正规子群（因而若 E 是有限维的，则 F(E) = GL(E)）。用 SL(E) 记由所有平延（习题 11）生成的 GL(E) 的正规子群；它包含在 F(E) 中。

\begin{enumerate}
\def\labelenumi{(\alph{enumi})}
\item
  若 \(\dim E \geqslant 2\) ，证明对 E 的每对非零向量 x, y，存在一个平延或两个平延的乘积，它把 x 映到 y（换句话说， \(\mathbf{SL}(\mathbf{E})\) 在 \(E = \{0\}\) 上传递地作用）。
\item
  设 V、W 是 E 的两个超平面， \(\dot{x}_{0} = x_{0} + V\) 是一个模 V 的剩余类，异于 V ，而 \(\dot{y}_{0} = y_{0} + W\) 是一个模 W 的剩余类，异于 W 。证明若 \(\dim E \geqslant 2\) ，则存在一个平延或两个平延的乘积，它把 V 映到 W 且把 \(\dot{x}_{0}\) 映到 \(\dot{y}_{0}\) （先考虑 V 与 W 相异的情形）。
\item
  若 \(\dim \mathbf{E} \geqslant 2\) ，证明任意两个异于恒等映射的平延在群 \(\mathbf{F}(\mathbf{E})\) 中共轭。
\item
  若 \(\dim \mathbf{E} \geqslant 3\) ，证明任意两个异于恒等映射的平延在群 \(\mathbf{SL}(\mathbf{E})\) 中共轭（借助 (b) 把它归结为两个平延的超平面相同的情形，然后利用 (a)）。
\item
  假设 \(\dim E = 2\) 。为使任意两个平延都在 \(\mathbf{SL}(\mathbf{E})\) 中共轭，必要且充分条件是子群 \(\mathbf{Q}\) （ \(\mathbf{K}^*\) 的、由 \(\mathbf{K}^*\) 的元素的平方生成的子群）与 \(\mathbf{K}^*\) 相同。（若 \(u\) 是一个异于恒等映射的平延， \(a\) 是 \(\mathbf{E}\) 中一个不在 \(u\) 下不变的向量，而 \(b = u(a) - a\) ，证明对每个平延 \(u'\) —— 与 \(u\) 在 \(\mathbf{SL}(\mathbf{E})\) 中共轭 —— 有 \(u'(a) - a = a\lambda + b\mu\) ，其中 \(\mu = 0\) 或 \(\mu \in \mathbf{Q}\) ；为此利用如下事实：在群 \(\mathbf{G}\) 中每个乘积 \(sts^{-1}t\) 都是平方的乘积。）
\item
  假设 \(\dim E \geqslant 2\) 。为使两个伸缩 \(u, u'\) 满足存在 \(v \in \mathbf{SL}(\mathbf{E})\) 使 \(u' = vuv^{-1}\) ，必要且充分条件是 \(u\) 与 \(u'\) 的类（习题 11）相同（利用 \((b)\) ）。由此推出，若一个伸缩的类包含在 \(\mathbf{K}^*\) 的换位子群中，则该伸缩属于群 \(\mathbf{SL}(\mathbf{E})\) （注意 \((vuv^{-1})u^{-1} = v(uv^{-1}u^{-1}))\) 。
\end{enumerate}

¶ 13. 设 E 是维数 ≥2 的右向量空间，且 H₀ 是 E 的一个超平面。

\begin{enumerate}
\def\labelenumi{(\alph{enumi})}
\item
  证明每个自同构 \(u\) （属于 \(\mathbf{F}(\mathbf{E})\) ，习题 12）是 \(\mathbf{SL}(\mathbf{E})\) 的一个自同构与可能还有超平面 \(\mathbf{H}_0\) 的一个伸缩的乘积（对 \(1_{\mathbf{E}} - u\) 的核的余维数用归纳法，并利用习题 12(a) 和 12(f)）。
\item
  证明 \(\mathbf{SL}(\mathbf{E})\) 包含 \(\mathrm{F(E)}\) 的换位子群，且除 \(\mathbf{E}\) 是含 2 个元素的域上的二维空间这一情形外与该群相同。（为证明 \(\mathbf{SL}(\mathbf{E})\) 包含 \(\mathrm{F(E)}\) 的换位子群，利用 (a) 和习题 \(12(f)\) 。为看出除所指出的情形外 \(\mathbf{SL}(\mathbf{E})\) 包含在 \(\mathrm{F(E)}\) 的换位子群中，证明对每个从 \(\mathrm{F(E)}\) 到交换群的同态，平延的像是单位元，利用习题 \(11(b)\) 和 \(12(c)\) 。）
\item
  证明 \(\mathbf{SL}(\mathbf{E})\) 等于其换位子群，除非 \(\dim \mathbf{E} = 2\) 且 \(\mathbf{K}\) 有 2 或 3 个元素。（若 \(\dim \mathbf{E} \geqslant 3\) ，注意每个平延 \(u\) 都可以写成 \(vwv^{-1}w^{-1}\) ，其中 \(v\) 是一个平延而 \(w \in \mathbf{SL}(\mathbf{E})\) ，利用习题 12(d)。若 \(\dim \mathbf{E} = 2\) ，首先注意每个自同构 \(u\) —— 它关于 \(\mathbf{E}\) 的一组基的矩阵具有形式 \(\begin{pmatrix} \lambda & 0 \\ 0 & \lambda^{-1} \end{pmatrix}\) —— 是平延的乘积，论证如 (a)；然后考虑换位子 \(uvu^{-1}v^{-1}\) ，其中 \(v\) 是平延，它关于同一组基的矩阵为 \(\begin{pmatrix} 1 & 1 \\ 0 & 1 \end{pmatrix}\) 。）
\end{enumerate}

¶ 14. 设 E 是域 K 上维数 ≥2 的向量空间。

\begin{enumerate}
\def\labelenumi{(\alph{enumi})}
\item
  证明射影群 PGL(E) 与线性群 GL(E) 对其中心（同构于 K 的中心乘法群，参见 § 1，习题 26）的商群典范同构。
\item
  设 \(\mathbf{PSL}(\mathbf{E})\) 表示群 \(\mathbf{SL}(\mathbf{E})\) 在 \(\mathbf{PGL}(\mathbf{E})\) 中的典范像；它是 \(\mathbf{PGL}(\mathbf{E})\) 的一个正规子群，包含 \(\mathbf{PGL}(\mathbf{E})\) 的换位子群（当 \(\mathbf{E}\) 有限维时）。证明 \(\mathbf{PSL}(\mathbf{E})\) 是 \(\mathbf{P}(\mathbf{E})\) 的一个双传递（参见 I，§ 5，习题 14）置换群。
\item
  设 \(a \neq 0\) 是 \(\mathbf{E}\) 的一个元素，并设 \(e = \pi(a) \in \mathbf{P}(\mathbf{E})\) （用 § 9 第 6 号的记号）。证明形如 \(x \mapsto x + a\phi(x)\) （其中 \(\phi \in \mathbf{E}^*\) 满足 \(\phi(a) = 0\) ）的平延组成 \(\mathbf{SL}(\mathbf{E})\) 的一个交换子群；若 \(\Gamma_e\) 是该子群在 \(\mathbf{PSL}(\mathbf{E})\) 中的像，则 \(\Gamma_e\) 是子群 \(\Phi_e\) （ \(\mathbf{PSL}(\mathbf{E})\) 中使 \(e\) 不变的子群）的正规子群。还证明与 \(\Gamma_e\) 共轭的子群（在 \(\mathbf{PSL}(\mathbf{E})\) 中）的并生成 \(\mathbf{PSL}(\mathbf{E})\) 。
\item
  设 \(\Sigma\) 是一个集合 F 上的本原群（I，§ 5，习题 13），它等于其换位子群。假设存在一个元素 \(c \in \mathbf{F}\) ，使得子群 \(\Phi_c\) （ \(\Sigma\) 的、使 \(c\) 不变的子群）包含一个交换正规子群 \(\Gamma_c\) ，使得与 \(\Gamma_c\) 共轭的子群（在 \(\Sigma\) 中）的并生成 \(\Sigma\) 。证明在这些条件下 \(\Sigma\) 是单群。（设 \(\Delta\) 是 \(\Sigma\) 的一个异于单位元的正规子群；利用 I，§ 5，习题 17，证明 \(\Sigma = \Delta\) . \(\Phi_c = \Delta\) . \(\Gamma_c\) 并利用 \(\Gamma_c\) 交换这一事实，推出 \(\Sigma\) 中的每个换位子都包含在 \(\Delta\) 中。）
\item
  由 (c)、(d) 和习题 13(c) 推出，群 \(\mathbf{PSL}(\mathbf{E})\) 是单群，除非 \(\dim \mathbf{E} = 2\) 且 \(\mathbf{K}\) 有 2 个或 3 个元素。
\end{enumerate}

\(^{*}(f)\) 若 \(\dim E = n + 1\) 且 K 是含 q 个元素的域，证明 PGL(E) 的阶为

\[
(q ^ {n + 1} - 1) (q ^ {n + 1} - q) \dots (q ^ {n + 1} - q ^ {n - 1}) q ^ {n}
\]

（利用(a)、§9的习题9(a)，以及K必然是交换域这一事实（V，§11，习题14和VIII，§11，第1号，定理1））。*

\begin{enumerate}
\def\labelenumi{(\alph{enumi})}
\setcounter{enumi}{6}
\item
  证明：若 \(\dim \mathbf{E} = 2\) 且 \(\mathbf{K}\) 是含 2 个（相应地，3 个）元素的域，则 \(\mathbf{PSL}(\mathbf{E})\) 同构于对称群 \(\mathfrak{S}_3\) （相应地，交错群 \(\mathfrak{A}_4\) ）。（把 \(\mathbf{PSL}(\mathbf{E})\) 看作 \(\mathbf{P}(\mathbf{E})\) 的一个置换群，注意它包含每个对换（相应地，每个循环（ \(a b c\) ）），并利用 \(f\) 。）
\item
  证明除 (g) 所考虑的情形外， \(\mathbf{SL}(\mathbf{E})\) 的每个异于 \(\mathbf{SL}(\mathbf{E})\) 的正规子群都包含在 \(\mathbf{SL}(\mathbf{E})\) 的中心中（利用 (d) 以及每个平延都是 \(\mathbf{SL}(\mathbf{E})\) 中的一个换位子（习题 13(c)））。
\end{enumerate}

¶ 15. 设 A 是一个环；对每个整数 \(n \geqslant 1\) ，把 \(\mathbf{A}^{n}\) 典范地等同于 \(\mathbf{A}^{(\mathbf{N})}\) 的由指标 \(\geqslant n\) 的坐标为 0 的元素组成的子模，并用 \(\mathbf{A}_{n}^{\prime}\) 记补子模，它由指标 \(<n\) 的坐标为 0 的元素组成。每个自同态 \(u \in \mathbf{GL}_{n}(\mathbf{A})\) 等同于 \(\mathbf{A}^{(\mathbf{N})}\) 的这样一个自同构：它在 \(\mathbf{A}^{n}\) 上的限制是 \(u\) ，在 \(\mathbf{A}_{n}^{\prime}\) 上的限制是恒等映射，于是 \(\mathbf{GL}_{n}(\mathbf{A})\) 被等同于 \(\mathbf{GL}(\mathbf{A}^{(\mathbf{N})})\) 的一个子群；用 F 记 \(\mathbf{GL}(\mathbf{A}^{(\mathbf{N})})\) 的由诸 \(\mathbf{GL}_{n}(\mathbf{A})\) 的并集组成的子群，换言之，即除有限个外使 \(\mathbf{A}^{(\mathbf{N})}\) 的典范基的元素都不变的子群（参见习题 12）。

设 \(\mathbf{T}_n\) 是 \(\mathbf{GL}_n(\mathbf{A})\) 的由诸 \(\mathrm{B}_{ij}(\lambda)\) （对 \(0 \leqslant i < n\) , \(0 \leqslant j < n\) , \(i \neq j\) , \(\lambda \in \mathbf{A}\) ，第 13 号）生成的子群，并令 \(\mathbf{T}\) 为 \(\mathbf{F}\) 的由诸 \(\mathbf{T}_n\) 的并集组成的子群。

\begin{enumerate}
\def\labelenumi{(\alph{enumi})}
\item
  证明对 \(n \geqslant 3\) ， \(T_n\) 等于其换位子群。
\item
  设 \(a, b\) 是 \(A\) 的两个可逆元素。证明
\end{enumerate}

\[
\left( \begin{array}{c c} a b & 0 \\ 0 & 1 \end{array} \right) = P \left( \begin{array}{c c} a & 0 \\ 0 & b \end{array} \right) Q \quad \text { and } \quad \left( \begin{array}{c c} a & 0 \\ 0 & b \end{array} \right) = P ^ {\prime} \left( \begin{array}{c c} b & 0 \\ 0 & a \end{array} \right) Q ^ {\prime}
\]

其中 \(P, Q, P', Q'\) 属于 \(T_2\) 。由此推出，矩阵 \(\begin{pmatrix} aba^{-1}b^{-1} & 0 \\ 0 & 1\end{pmatrix}\) 属于 \(T_2\) 。

\begin{enumerate}
\def\labelenumi{(\alph{enumi})}
\setcounter{enumi}{2}
\tightlist
\item
  由 (b) 推出， \(\mathbf{GL}_{n}(\mathbf{A})\) 的换位子群包含在 \(T_{2n}\) 中（在 (b) 中把 A 换成 \(\mathbf{M}_{n}(\mathbf{A})\) ）。（利用 (a)）得出结论：F 的换位子群等于 T。
\end{enumerate}

\begin{enumerate}
\def\labelenumi{\arabic{enumi}.}
\setcounter{enumi}{15}
\tightlist
\item
  \begin{enumerate}
  \def\labelenumii{(\alph{enumii})}
  \tightlist
  \item
    设 \(U\) 是一个 \(n\) 阶方阵（在域 \(\mathbf{Q}\) 上），其所有非对角元素都等于同一个有理数 \(r > 0\) ，且其对角线 \((d_1, \ldots, d_n)\) 满足： \(d_i \geqslant r\) （对所有 \(i\) ）且 \(d_i > r\) （除 \(i\) 的可能一个值外）。证明 \(U\) 可逆（若 \(x = (x_i)_{1 \leqslant i \leqslant n}\) 是方程 \(U.x = 0\) 的一个解，证明这些 \(x_i\) 都有相同的符号，并由此推出它们全为零）。
  \end{enumerate}
\end{enumerate}

\begin{enumerate}
\def\labelenumi{(\alph{enumi})}
\setcounter{enumi}{1}
\tightlist
\item
  设 E 是含 m 个元素的有限集，这些元素记作 \(a_{i}\) ( \(1 \leqslant i \leqslant m\) )，而 \((\mathbf{A}_{j})_{1 \leqslant j \leqslant n}\) 是 E 的 n 个子集所成的族，彼此互不相同；假设 \(\text{Card}(\mathbf{A}_{i} \cap \mathbf{A}_{j}) = r \geqslant 1\) 对每个有序对 \((i, j)\) （由 \([1, n]\) 的不同元素组成的有序对）成立。证明必然有 \(n \leqslant m\) 。（令 \(V = (c_{ij})\) 是 \((m, n)\) 型矩阵，使得 \(c_{ij} = 1\) （当 \(a_{i} \in A_{j}\) 时），否则 \(c_{ij} = 0\) 。把 (a) 应用于矩阵 \(^{t}V\) . V （n 阶）。）
\end{enumerate}

¶ 17. 设 K 为一个域（交换或非交换），E 为 K 上的 n 维右向量空间，u 为 E 的一个自同态。

\begin{enumerate}
\def\labelenumi{(\alph{enumi})}
\item
  证明：若 \(n > 1\) 且 \(u\) 不是中心位似，则存在 \(E\) 的一组基，使得 \(u\) 关于这组基的矩阵具有形式 \((\alpha_{ij})\) ，其中 \(\alpha_{ij} = 0\) （对 \(j > i + 1\) ）且 \(\sum_{i=1}^{n-1} \alpha_{i,i+1} = 1\) （对 \(n\) 用归纳法论证）。
\item
  证明：若 \(n > 2\) ，则存在 \(\mathbf{E}\) 的一组基，使得 \(u\) 关于这组基的矩阵具有形式 \((\alpha_{ij})\) ，其中 \(\alpha_{ij} = 0\) （对 \(j > i + 1\) ）且
\end{enumerate}

\(\sum_{i=1}^{n-1} \alpha_{i, i+1} = 0\) 。（对 \(n\) 用归纳法并利用 \((a)\) 论证，注意对 \(u\) 或对 \(u + \gamma 1_{E}\) （其中 \(\gamma\) 属于 \(K\) 的中心）证明该命题是等价的，并分别研究 \(u\) 的秩为 1 或 2 的情形。）

\begin{enumerate}
\def\labelenumi{(\alph{enumi})}
\setcounter{enumi}{2}
\item
  设 \(A = (\alpha_{ij})\) 是具有 (b) 中所述性质的一个矩阵。又设 \(S\) 是阶为 \(n(\varepsilon_{ij})\) 的矩阵，使得 \(\varepsilon_{ij} = 0\) ，除非 \(i = j + 1\) ，这时 \(\varepsilon_{j+1,j} = 1\) 。证明存在一个矩阵 \(B = (\beta_{ij})\) （ \(n\) 阶），使得 \(A = \lambda E_{nn} + (BS - SB)\) （ \(\lambda \in \mathbf{K}\) ）。（取 \(B\) 使得 \(\beta_{i1} = 0\) （对 \(1 \leqslant i \leqslant n\) ）且 \(\beta_{ij} = 0\) （对 \(j > i + 2\) ）。）
\item
  假设 K 是交换的。由 (c) 推出，E 的每个满足 \(\operatorname{Tr}(u) = 0\) 的自同态 u 都可以写成 vw - wv，其中 v 和 w 是 E 的自同态。
\end{enumerate}

